% !TEX root = ../main.tex
% ---------------------------------------------------------------- хураангуй
\chapter*{Хураангуй}
\addcontentsline{toc}{chapter}{Хураангуй}

\textbf{Асуудал.} Олон цөмтэй компьютерыг ашиглахын тулд ажлыг thread-үүдэд хуваах шаардлагатай.
Гэхдээ хуваасан ажлын хэсгүүд хэн ямар өгөгдөлд бичих, хэзээ дуусахыг хэрхэн хүлээх, үр дүнг хэрхэн
нэгтгэхийг зөв зохицуулахгүй бол үр дүн буруу гарна. \emph{Operating System Concepts} (10-р хэвлэл) номын
4-р бүлгийн “Multithreaded Sorting Application” төсөл энэ асуудлыг эрэмбэлэлтийн жишээн дээр судалдаг.

\textbf{Шийдэл.} Java 17 дээр жагсаалтыг хоёр хуваан, хоёр \emph{sorting thread} тус бүр параметрээр өгсөн
мужаа эрэмбэлж, гурав дахь \emph{merging thread} эх массивтай ижил хэмжээтэй хоёр дахь global массив руу
нэгтгэдэг програм хийв. Parent thread бүх thread-ийг \code{join()} хийсний дараа үр дүнг хэвлэнэ. Синхрончлолыг
бүхэлд нь \code{start()}/\code{join()} хангадаг тул lock шаардлагагүй. Таван алгоритм (insertion, bubble,
selection, merge, quick), оролтын шалгалт, JSON event trace хийсэн. Мөн бодит thread-тэй холбогдсон хөтчийн
\emph{Sorting animation} хийсэн бөгөөд хоёр горимтой: бүртгэсэн алхмыг тоглуулах \emph{Replay}, амьд thread-ийг
\code{Semaphore}-оор нэг үйлдлээр урагшлуулах \emph{Lockstep}.

\textbf{Туршилт.} JUnit-гүй 100 автомат тест бичиж, macOS (Temurin 17) ба Linux (OpenJDK 21) орчинд ажиллуулав.
Тест нь номын жишээ, захын тохиолдол, буруу оролт, давхардал, $n = 10^6$ хүртэлх том өгөгдөл, thread-ийн
дараалал, trace, Lockstep-ийг хамарна. Хугацааг 1 thread, 1 thread хувааж, 3 thread гэсэн гурван хувилбараар,
3 тусдаа JVM дээр хэмжив.

\textbf{Гол үр дүн.} 100/100 тест амжилттай. Sorter бүр зөвхөн өөрийн мужид бичдэгийг trace-ээс баталсан.
Хоёр цөмтэй машин дээр merge ба quick sort-ын хурдсалт $n = 10^6$ үед $1.7\times$ орчим гарсан. Энэ нь 2-оос
бага бөгөөд шалтгаан нь нэгтгэх үе дан thread-ээр явдагт оршино (Амдалын хууль). O($n^2$) алгоритмд
($n = 20\,000$) insertion ба selection sort-ын хурдсалт $2.4$--$3.6\times$ гарсан ч ашгийн ихэнх нь массивыг
хуваасантай холбоотой, учир нь хоёр хэсгийг эрэмбэлэх ажил $2(n/2)^2 = n^2/2$ байна. $n \le 1000$ үед thread
үүсгэх, \code{join()} хийх зардал давамгайлж, олон урсгалт хувилбар нэг урсгалтаас удаан ажилласан.

\medskip
\noindent\textbf{Түлхүүр үг:} thread, олон урсгалт програмчлал, \code{join}, зэрэгцээ тооцоолол,
эрэмбэлэлт, merge, semaphore, Амдалын хууль, Java.
